% 火星（综述）
% license CCBYSA3
% type Wiki

本文根据 CC-BY-SA 协议转载翻译自维基百科\href{https://en.wikipedia.org/wiki/Mars}{相关文章}。

\begin{figure}[ht]
\centering
\includegraphics[width=6cm]{./figures/b86f152bd3571454.png}
\caption{火星的真实颜色\(^{\text{a}}\)，由“希望号”轨道器拍摄。画面中央可见塔尔西斯山脉，左侧为奥林匹斯山，右侧则是水手谷。} \label{fig_Mars_1}
\end{figure}
火星是距离太阳的第四颗行星。由于其呈橙红色的外观，它也被称为“红色星球”\(^{\text{[22,23]}}\)。火星是一颗类似沙漠的岩质行星，大气稀薄，主要成分为二氧化碳（CO\(_2\)）。在平均地表高度处，其大气压仅为地球的几千分之一，大气温度范围为 −153 至 20 °C（−243 至 68 °F）\(^{\text{[24]}}\)，宇宙辐射水平很高。火星保留了一部分水，存在于地下以及稀薄大气中，可形成卷云、雾气、霜冻、较大的极地永久冻土区与冰帽（伴随季节性的 CO\(_2\) 雪），但其表面不存在液态水体。其表面重力约为地球的三分之一，为月球的两倍。直径为 6,779 公里（4,212 英里），约为地球的一半，是月球的两倍，其表面积则相当于地球全部陆地面积的总和。

细微尘埃遍布火星表面和大气，在低重力条件下，即便是稀薄大气的微弱风力也能将其扬起并扩散开来。火星地形大致呈现显著的南北分界，即“火星二元性结构”：北半球主要是平坦、低洼的平原，而南半球则由遍布陨石坑的高地组成。从地质学上看，火星处于相对活跃的状态，地下会发生“火震”，但其上也分布着许多巨型而已灭绝的火山（最高的是奥林匹斯山，高 21.9 公里或 13.6 英里），以及太阳系中最大峡谷之一的水手谷（长约 4,000 公里或 2,500 英里）。火星有两颗天然卫星，形状小而不规则：火卫一（Phobos）与火卫二（Deimos）。由于其轴倾角为 25 度，火星像地球（轴倾角 23.5 度）一样经历季节变化。一个火星太阳年相当于 1.88 个地球年（687 个地球日），一个火星太阳日（sol）则等于 24.6 小时。

火星与其他行星一同形成于大约 45 亿年前。在火星的诺亚纪（Noachian period，约 45 亿年至 35 亿年前）期间，其地表经历了陨石撞击、峡谷形成、侵蚀作用、可能存在的海洋，以及磁层的消失。赫斯珀里亚纪（Hesperian period，自约 35 亿年前开始，结束于约 33–29 亿年前）以广泛的火山活动和导致巨大泄流河道形成的洪水为主要特征。亚马逊纪（Amazonian period）从其后延续至今，是当前地质过程的主要影响因素。由于火星的地质历史，对于火星过去或当前是否存在生命的探索仍是活跃的科学研究领域，其中一些可能的生命迹象仍需进一步确认。

由于在地球夜空中可被肉眼看到、呈红色并缓慢移动，火星自古以来就被人类观测，并在不同文化中具有多样的象征意义。1963 年，人类首次发射前往火星的探测器“火星 1 号”（Mars 1），但途中失联。第一项成功飞越火星的探测任务发生在 1965 年，由“水手 4 号”（Mariner 4）完成。1971 年，“水手 9 号”（Mariner 9）进入火星轨道，成为人类历史上首个进入除月球、太阳或地球以外天体轨道的航天器；同年，“火星 2 号”成为首个在火星发生非受控撞击的探测器，“火星 3 号”则实现了首次成功着陆。自 1997 年以来，探测器在火星上持续运行。有时，超过十个探测器会同时在火星轨道或地表开展活动，这一数量超过地球以外任意行星。

火星常被视为未来载人探索任务的目标，尽管目前尚无正式计划付诸实施。
\subsection{自然历史}  
\subsubsection{形成} 
科学家推测，在太阳系形成过程中，火星是由于围绕太阳运行的原行星盘中的物质通过随机的失控聚积过程形成的。火星在太阳系中的位置使其具备许多独特的化学特征。相较于地球，沸点较低的元素，例如氯、磷和硫，在火星上要常见得多；这些元素很可能是在年轻太阳强烈的太阳风作用下被推移到更外侧区域\(^{\text{[25]}}\)。

晚期重轰炸  
在行星形成之后，内太阳系可能经历了所谓的“晚期重轰炸”。火星约有 60\% 的表面保留了那个时期的撞击记录\(^{\text{[26,27,28]}}\)，而其余许多区域的地表之下，可能也存在由同一时期事件形成的巨大撞击盆地。然而，最近的建模研究对晚期重轰炸的存在提出了质疑\(^{\text{[29]}}\)。证据显示，火星北半球可能存在一个巨大的撞击盆地，范围达 10,600 × 8,500 公里（6,600 × 5,300 英里），约为月球南极—艾特肯盆地面积的四倍；若得到确认，这将是迄今发现的最大撞击盆地\(^{\text{[30]}}\)。科学家推测该盆地形成于约 40 亿年前，当时火星被一个与冥王星大小相当的天体撞击，造成了火星南北半球的显著地形差异，并形成覆盖火星 40\% 表面的平滑的北方平原（Borealis basin）\(^{\text{[31,32]}}\)。

一项 2023 年的研究根据火卫二（Deimos）的轨道倾角提出证据，显示火星在 35 亿至 40 亿年前可能曾拥有一个环系统\(^{\text{[33]}}\)。该环系统可能由一颗质量约为火卫一（Phobos）20 倍的古老卫星解体形成；而火卫一可能正是这一原始环系统的残余物\(^{\text{[34,35]}}\)。
\begin{figure}[ht]
\centering
\includegraphics[width=14.25cm]{./figures/4b6ac302fa80c765.png}
\caption{} \label{fig_Mars_2}
\end{figure}
\begin{figure}[ht]
\centering
\includegraphics[width=14.25cm]{./figures/7861217b62cae434.png}
\caption{} \label{fig_Mars_3}
\end{figure}
火星的地质历史可以被划分为多个时期，但以下三个时期最为主要 \(^{\text{[36,37]}}\)。
\begin{itemize}
\item 诺亚纪（Noachian period）：约 45 亿年至 35 亿年前，是火星现存最古老地表的形成阶段。诺亚纪地表布满大量大型撞击坑。塔尔西斯隆起（Tharsis bulge）这一火山高地被认为形成于该时期，其末期可能经历了大规模液态水泛滥。该时期名称源自 Noachis Terra \(^{\text{[38]}}\)。
\item 赫斯珀纪（Hesperian period）：约 35 亿年至介于 33 亿年至 29 亿年前之间。赫斯珀纪的特征是广泛的熔岩平原的形成。名称源自 Hesperia Planum \(^{\text{[38]}}\)。
\item 亚马孙纪（Amazonian period）：约 33 亿年至 29 亿年前至今。亚马孙纪区域的陨石撞击坑数量较少，但地貌表现多样。奥林匹斯山（Olympus Mons）形成于该时期，同时火星其他地区也存在熔岩流活动。名称源自 Amazonis Planitia \(^{\text{[38]}}\)。
\end{itemize}
\subsubsection{近期地质活动}
火星上仍然存在地质活动。阿萨巴斯卡峡谷（Athabasca Valles）中有形成于约 2 亿年前的片状熔岩流。刻耳柏洛斯裂谷（Cerberus Fossae）这一断陷带中的水流活动发生于不足 2,000 万年前，显示出同样年轻的火山侵入事件 \(^{\text{[39]}}\)。火星勘测轨道器（Mars Reconnaissance Orbiter）已拍摄到雪崩的影像 \(^{\text{[40,41]}}\)。
\subsection{物理特征}
\begin{figure}[ht]
\centering
\includegraphics[width=8cm]{./figures/23fb842f59355847.png}
\caption{按比例绘制的火星及内太阳系中具有行星质量的天体。从左至右依次为：水星、金星、地球、月球、火星和谷神星。} \label{fig_Mars_4}
\end{figure}
火星的直径约为地球的一半，或约为月球的两倍，其表面积仅略小于地球全部陆地面积 \(^{\text{[2]}}\)。火星的密度低于地球，体积约为地球的 15\%，质量约为地球的 11\%，因此其表面重力约为地球的 38\%。火星是目前已知的唯一一颗沙漠行星，即表面类似于地球荒漠环境的类地岩质行星。火星表面呈红橙色，是由于铁(III)氧化物（纳米级 Fe\(_2\)O\(_3\)）以及铁(III)氧化物-氢氧化物矿物针铁矿（goethite）所致 \(^{\text{[42]}}\)。其颜色有时类似奶油太妃糖 \(^{\text{[43]}}\)；其他常见的地表颜色还包括金色、棕色、褐色以及因矿物组成不同而呈现的绿ish色调 \(^{\text{[43]}}\)。
\subsubsection{内部结构}
\begin{figure}[ht]
\centering
\includegraphics[width=8cm]{./figures/4694e492a3c0d60d.png}
\caption{火星的内部结构 \(^{\text{[44,45]}}\)} \label{fig_Mars_5}
\end{figure}
与地球类似，火星也分异为高密度的金属核以及其上覆盖的低密度岩质层 \(^{\text{[46][47]}}\)。最外层为地壳，平均厚度约为 42–56 公里（26–35 英里） \(^{\text{[48]}}\)，在伊西迪斯平原（Isidis Planitia）处最薄，为 6 公里（3.7 英里），在塔尔西斯高原（Tharsis plateau）南部最厚，可达 117 公里（73 英里） \(^{\text{[49]}}\)。相比之下，地球地壳平均厚度为 27.3 ± 4.8 公里 \(^{\text{[50]}}\)。火星地壳中最丰富的元素包括硅、氧、铁、镁、铝、钙和钾。火星已被确认具有地震活动 \(^{\text{[51]}}\)；在 2019 年，报道指出 InSight 着陆器已探测并记录超过 450 次火震及相关事件 \(^{\text{[52][53]}}\)。

在地壳之下是硅酸盐地幔，它形成了火星表面许多构造与火山特征。火星上地幔的上部为低速带，在此区段中地震波速度低于邻近深度层。地幔似乎在约 250 公里深度范围内保持刚性 \(^{\text{[45]}}\)，使得火星的岩石圈相较地球更为厚实。在此深度以下，地幔逐渐变得更具延展性，且地震波速度再次开始增加 \(^{\text{[54]}}\)。火星地幔似乎并不存在类似地球下地幔的热绝缘层；相反，在约 1050 公里深度以下，其矿物学特征与地球的过渡带相似 \(^{\text{[44]}}\)。在地幔底部存在一层约 150–180 公里厚的基性液态硅酸盐层 \(^{\text{[45][55]}}\)。火星地幔看似高度非均质，其中包含直径可达 4 公里的高密度碎块，可能是在约 45 亿年前的巨型撞击中被深度注入；来自八次火震的高频波在穿越这些局部区域时出现减速，模型显示这些非均质体是组成上独特的残余碎屑，由于火星缺乏板块构造且内部对流迟缓，无法实现完全均质化，因此得以保存 \(^{\text{[56][57]}}\)。

火星的铁镍核至少部分为熔融状态，并可能具有固体内核 \(^{\text{[58][59][60]}}\)。其半径约为火星半径的一半，大约为 1650–1675 公里，且富含硫、氧、碳和氢等轻元素 \(^{\text{[61][62]}}\)。其核心温度估计为 2000–2400 K \(^{\text{[63]}}\)，相比之下，地球固体内核的温度为 5400–6230 K。2025 年，基于 InSight 着陆器的数据，一组研究人员报告探测到一个半径为 613 公里（381 英里） ± 67 公里（42 英里）的固体内核 \(^{\text{[60]}}\)。
\subsubsection{表面地质}  
更多信息参见：火星风化层（Martian regolith）。火星是一颗类地行星，其表面由含硅和氧的矿物、金属以及其他通常构成岩石的元素组成。火星表面主要由拉斑玄武岩（tholeiitic basalt）构成 \(^{\text{[64]}}\)，但部分区域的含硅量高于典型玄武岩，可能类似地球的安山岩或硅质玻璃。低反照率区域显示斜长石长石的富集，而北部低反照率地区呈现较高含量的层状硅酸盐与高硅玻璃。南部高地的部分地区可检测到高钙辉石的存在。此外，还发现局地富集的赤铁矿与橄榄石 \(^{\text{[65]}}\)。火星大部分表面深覆以细粒的铁(III)氧化物尘埃 \(^{\text{[66]}}\)。

凤凰号着陆器（Phoenix lander）返回的数据表明火星土壤略呈碱性，且含有镁、钠、钾及氯等元素。这些养分同样存在于地球土壤中，并对植物生长必不可少 \(^{\text{[67]}}\)。着陆器的实验显示火星土壤的 pH 为 7.7，为弱碱性，并按重量含有 0.6\% 的高氯酸盐 \(^{\text{[68][69]}}\)，这种浓度对人类具有毒性 \(^{\text{[70][71]}}\)。

条纹（streaks）在火星上十分常见，并且在陨石坑、槽谷与峡谷的陡坡上频繁出现新的条纹。这些条纹初期颜色较深，随后随时间变浅。条纹可从极小区域开始，并扩展数百米，其路径会绕过巨石等障碍物。最普遍接受的假说认为，这些条纹是由明亮尘埃的雪崩或尘卷风作用暴露出的暗色下层土壤 \(^{\text{[72]}}\)。其他提出的解释包括与水有关的机制，甚至可能涉及生物生长的假说 \(^{\text{[73][74]}}\)。

火星表面的环境辐射水平平均为每日 0.64 毫西弗，显著低于往返火星途中每日 1.84 毫西弗（或每日 22 毫拉德）的辐射水平 \(^{\text{[75][76]}}\)。相比之下，地球近地轨道（低地轨道）中，例如空间站所在轨道的辐射约为每日 0.5 毫西弗 \(^{\text{[77]}}\)。赫拉斯平原（Hellas Planitia）具有火星表面最低的辐射水平，约每日 0.342 毫西弗，其位于哈德里亚卡斯山（Hadriacus Mons）西南方的熔岩管（lava tubes）中辐射水平可能低至每日 0.064 毫西弗 \(^{\text{[78]}}\)，与地球飞机飞行期间的辐射水平相当。
\subsubsection{磁学特征} 
虽然火星不存在有组织的全球磁场的证据 \(^{\text{[79]}}\)，但观测表明其地壳的部分区域已被磁化，这暗示其偶极磁场在过去曾发生过极性反转。磁化敏感矿物所记录的这种古地磁性与地球洋底发现的交替磁带类似。一项发表于 1999 年并于 2005 年 10 月利用火星全球探测者（Mars Global Surveyor）重新分析的假说提出，这些磁带可能表明火星在约 40 亿年前曾具有板块构造活动，当时行星发电机仍在运作，而全球磁场尚未衰亡 \(^{\text{[80]}}\)。
\subsection{地理与地貌}
\begin{figure}[ht]
\centering
\includegraphics[width=8cm]{./figures/6c02165237502845.png}
\caption{标注了各地貌特征的火星地形图，并清晰显示出火星二分性（北部低洼半球与南部高地半球）。} \label{fig_Mars_6}
\end{figure}
尽管更以绘制月球地图而闻名，约翰·海因里希·冯·梅德勒（Johann Heinrich von Mädler）与威廉·贝尔（Wilhelm Beer）是最早的火星地理绘图者。他们首先确认火星的大部分地表特征是永久性的，并更精确地测定了火星的自转周期。1840 年，梅德勒将十年的观测资料整合，绘制出第一幅火星地图 \(^{\text{[81]}}\)。

火星表面的特征名称来源多样。反照率特征以古典神话命名；直径大于约 50 公里的陨石坑以已故科学家、作家及其他对火星研究有贡献的人命名；较小的陨石坑则以全球人口不足 10 万的城镇命名。大型峡谷以不同语言中“火星”或“星”的词汇命名；较小的峡谷则以河流命名 \(^{\text{[82]}}\)。

大型反照率特征保留了许多古老的命名，但常根据对其性质的新认识进行更新。例如，原名为 Nix Olympica（奥林匹斯之雪）的特征现名为 Olympus Mons（奥林匹斯山）\(^{\text{[83]}}\)。从地球上观测，火星表面按反照率差异分为两类区域。较浅色的平原被红铁氧化物丰富的尘埃与沙覆盖，曾被认为是火星的“大陆”，因而被命名为阿拉伯高地（Arabia Terra，阿拉伯之地）或亚马逊平原（Amazonis Planitia，亚马逊平原）。暗色特征则被认为是海洋，因此命名包括 Mare Erythraeum、Mare Sirenum 与 Aurorae Sinus。地球上可见的最大暗区是 Syrtis Major Planum \(^{\text{[84]}}\)。永久的北极冰盖称为 Planum Boreum，南极冰盖称为 Planum Australe \(^{\text{[85]}}\)。

火星赤道由其自转定义，但其本初子午线的位置与地球相同，由人为选定：梅德勒和贝尔在 1830 年制作的第一幅火星地图上选定了一条经线。1972 年“海盗 9 号”（Mariner 9）提供了大量火星影像后，位于子午线海（Sinus Meridiani，“中央海湾”或“子午线海湾”）中的一个小陨石坑（后称 Airy-0）被梅尔顿·E·戴维斯、哈罗德·马苏斯基与热拉尔·德·沃科勒选为 0.0° 经度的标定点，以保持与最初选定的位置一致 \(^{\text{[86][87][88]}}\)。

由于火星没有海洋，因此也不存在“海平面”；必须选取一个零高度面作为参考，即火星的准地球体面（areoid），它与地球的地球体（geoid）相对应 \(^{\text{[89][90]}}\)。零高度被定义为大气压力为 610.5 Pa（6.105 mbar）的高度 \(^{\text{[91]}}\)。该压力对应于水的三相点，大约是地球海平面气压的 0.6\%（0.006 atm）\(^{\text{[92]}}\)。

为制图目的，美国地质调查局将火星表面划分为三十个测图分区，每个分区以其中包含的经典反照率特征命名 \(^{\text{[93]}}\)。2023 年 4 月，《纽约时报》报道了一幅基于“希望号”（Hope）探测器影像更新的火星全球地图 \(^{\text{[94]}}\)。另一幅相关但细节更高的全球火星地图由 NASA 于 2023 年 4 月 16 日发布 \(^{\text{[95]}}\)。
\subsubsection{火山}
\begin{figure}[ht]
\centering
\includegraphics[width=8cm]{./figures/cc4c46ae5220f270.png}
\caption{火星最高的火山——奥林匹斯山（Olympus Mons）的图像。其横跨范围约为 550 公里（340 英里）。} \label{fig_Mars_7}
\end{figure}
广阔的塔尔西斯（Tharsis）高原地区包含若干巨型火山，其中包括盾状火山奥林匹斯山（Olympus Mons）。其火山体宽度超过 600 公里（370 英里）\(^{\text{[96][97]}}\)。由于这座山体规模极其庞大，且其边缘具有复杂结构，因此难以给出确切的高度。从其西北边缘峭壁的山脚至山顶的局部高差超过 21 公里（13 英里）\(^{\text{[97]}}\)，略高于从海底基部测量的冒纳凯阿火山（Mauna Kea）高度的两倍。从亚马孙平原（Amazonis Planitia）出发，向西北方向跨越超过 1,000 公里（620 英里）至火山峰顶的整体海拔变化接近 26 公里（16 英里）\(^{\text{[98]}}\)，约为珠穆朗玛峰（Mount Everest，海拔略高于 8.8 公里〔5.5 英里〕）的三倍。因此，奥林匹斯山是太阳系中最高或第二高的山体；唯一可能更高的是灶神星（Vesta）上的雷亚西尔维亚（Rheasilvia）中央峰，高度约为 20–25 公里（12–16 英里）\(^{\text{[99]}}\)。
\subsubsection{撞击地形}  
火星地形的二元性极为显著：被熔岩流抹平的北部平原，与因古老撞击而密布坑洼的南部高地形成强烈对比。可能在约 40 亿年前，火星北半球曾被一个大小约为地球月球十分之一至三分之二的天体撞击。如果这一假说成立，则火星北半球将是一个尺寸达 10,600 × 8,500 公里（6,600 × 5,300 英里）的撞击盆地，其面积大致相当于欧洲、亚洲与澳大利亚总和，超过乌托邦平原（Utopia Planitia）以及月球南极—艾特肯盆地（South Pole–Aitken Basin），成为太阳系中最大的撞击盆地 \(^{\text{[100][101][102]}}\)。

火星上直径至少 5 公里（3.1 英里）的撞击坑数量达 43,000 个 \(^{\text{[103]}}\)。暴露最明显的大型撞击坑为赫拉斯盆地（Hellas），其直径 2,300 公里（1,400 英里），深达 7,000 米（23,000 英尺），并作为一个明亮的反照率特征从地球上即可清晰观测到 \(^{\text{[104][105]}}\)。其他重要的撞击构造包括阿耳戈里平原（Argyre），其直径约为 1,800 公里（1,100 英里）\(^{\text{[106]}}\)，以及伊西迪斯盆地（Isidis），其直径约为 1,500 公里（930 英里）\(^{\text{[107]}}\)。由于火星的质量和尺寸较小，其受到天体撞击的概率约为地球的一半；但由于火星更靠近小行星带，它受到来自该区域物质撞击的几率更高。火星也更容易遭到短周期彗星的撞击，即那些轨道位于木星轨道以内的彗星 \(^{\text{[108]}}\)。

部分火星撞击坑呈现的形态特征表明地表可能在撞击后变得潮湿 \(^{\text{[109]}}\)。
\subsubsection{构造地点}
\begin{figure}[ht]
\centering
\includegraphics[width=8cm]{./figures/dabff5c6b996abd3.png}
\caption{由“海盗 1 号”（Viking 1）探测器拍摄的水手谷（Valles Marineris）。} \label{fig_Mars_8}
\end{figure}
巨型峡谷水手谷（Valles Marineris，其拉丁语意为“水手号峡谷”，在早期运河地图中亦称为 Agathodaemon \(^{\text{[110]}}\)）全长约 4,000 公里（2,500 英里），深度可达 7 公里（4.3 英里）。其长度相当于欧洲的全长，并横跨火星周长的五分之一。相比之下，地球上的科罗拉多大峡谷（Grand Canyon）仅长 446 公里（277 英里），深近 2 公里（1.2 英里）。水手谷的形成被认为源于塔尔西斯（Tharsis）地区的隆起，使得水手谷区域的地壳发生断陷。2012 年，一项研究提出水手谷不仅仅是一个地堑结构，而可能是一条板块边界，其上发生了约 150 公里（93 英里）的横向位移，这意味着火星可能拥有由两块构成的板块构造体系 \(^{\text{[111][112]}}\)。
\subsubsection{洞穴与孔洞} 
美国宇航局“火星奥德赛号”（Mars Odyssey）轨道器搭载的热发射成像系统（THEMIS）拍摄的图像在阿耳狄西亚山（Arsia Mons）的山坡上揭示了七个可能的洞穴入口 \(^{\text{[113]}}\)。这些洞穴以发现者亲人的名字命名，被统称为“七姐妹”（seven sisters）\(^{\text{[114]}}\)。洞穴入口宽度从 100 到 252 米（328 至 827 英尺）不等，推测深度至少为 73 至 96 米（240 至 315 英尺）。由于大多数洞穴的底部无法被光照及成像到，它们可能远比这些估计值更深，并在地下进一步扩展。“Dena”是唯一例外；其底部可见，并测得深度为 130 米（430 英尺）。这些洞穴内部可能免受微陨石、紫外辐射、太阳耀斑及轰击火星表面的高能粒子的影响，从而形成受保护的环境 \(^{\text{[115][116]}}\)。
\subsubsection{其他特征}
\begin{figure}[ht]
\centering
\includegraphics[width=6cm]{./figures/6aecb1f42844b25f.png}
\caption{艺术家概念图，展示从火星间歇泉喷出的夹沙喷流（由 NASA 发布，艺术家：Ron Miller）。} \label{fig_Mars_9}
\end{figure}
火星间歇泉（Martian geysers，或称二氧化碳喷口 CO\(_2\) jets）被认为是火星南极地区在春季融化期间发生的小规模气体与尘埃喷发地点。“暗沙丘斑点”（dark dune spots）与“蜘蛛”或蜘蛛状地貌（araneiforms）是这类喷发所导致的两种最显著的地表特征 \(^{\text{[117]}}\)。

本节摘自“火星风化层（Martian regolith）”中“§ 大气尘埃（Atmospheric dust）”部分。

轨道上观测到的火星尘暴细节。

火星在没有尘暴时（左图，2001 年 6 月）与发生全球性尘暴时（右图，2001 年 7 月）的比较图像，由火星全球勘测者（Mars Global Surveyor）拍摄。

2018 年 5 月至 9 月期间的火星尘暴光学深度（tau）变化，由火星气候探测仪（Mars Climate Sounder）测量。

尘埃与水云的差异：图像下部中央的黄色云层为大型尘埃云，而其他白色云层为水云。

与地球相比，类似粒径的尘埃会更快从较稀薄的火星大气中沉降。例如，2001 年火星全球性尘暴扬起的尘埃仅在大气中停留约 0.6 年，而皮纳图博火山（Mount Pinatubo）喷发的尘埃在地球大气中约需两年才完全沉降 \(^{\text{[118]}}\)。然而，在当前火星环境条件下，涉及的物质迁移量一般远小于地球。即便是 2001 年火星全球尘暴，其移动的尘埃量折算成均匀沉积层也仅约 3 微米厚，分布在南纬 58°至北纬 58°之间 \(^{\text{[118]}}\)。在两处火星车着陆点的尘埃沉积速度大约为每 100 个火星日（sol）沉积一粒尘埃颗粒的厚度 \(^{\text{[119]}}\)。
\begin{figure}[ht]
\centering
\includegraphics[width=8cm]{./figures/9be5f19566825e95.png}
\caption{暗沙丘斑点} \label{fig_Mars_10}
\end{figure}\begin{figure}[ht]
\centering
\includegraphics[width=8cm]{./figures/9e6fe720e071c2c8.png}
\caption{从轨道上观测到的火星尘暴细节。} \label{fig_Mars_11}
\end{figure}\begin{figure}[ht]
\centering
\includegraphics[width=8cm]{./figures/200ce86596732349.png}
\caption{2001 年 6 月无尘暴的火星（左）与 2001 年 7 月发生全球性尘暴的火星（右），由火星全球勘测者（Mars Global Surveyor）拍摄。} \label{fig_Mars_12}
\end{figure}\begin{figure}[ht]
\centering
\includegraphics[width=8cm]{./figures/35c9fae96a3ee78b.png}
\caption{尘埃云与水云的差异：图像下部中央的黄色云层是大型尘埃云，其他白色云层则是水云。} \label{fig_Mars_13}
\end{figure}
\begin{figure}[ht]
\centering
\includegraphics[width=10cm]{./figures/7e0122ca60745c31.png}
\caption{由“好奇号”（Curiosity）火星车拍摄的火星尘卷风。/一个直径 30 米、高度 800 米的尘卷风。曾观测到高度达 20 公里的尘卷风。/尘卷风在火星表面形成扭曲的深色轨迹。 } \label{fig_Mars_14}
\end{figure}
\subsection{大气}
\begin{figure}[ht]
\centering
\includegraphics[width=8cm]{./figures/f6a34681d67508ea.png}
\caption{} \label{fig_Mars_15}
\end{figure}
火星在约 40 亿年前失去了其磁层 \(^{\text{[120]}}\)，可能是由于多次小行星撞击所致 \(^{\text{[121]}}\)，因此太阳风可直接与火星电离层相互作用，通过剥离外层大气原子而降低大气密度 \(^{\text{[122]}}\)。火星全球勘测者（Mars Global Surveyor）与火星快车号（Mars Express）均探测到火星背后延伸入太空的电离大气粒子尾迹 \(^{\text{[120][123]}}\)，这种大气逃逸现象正由 MAVEN 轨道器进行研究。与地球相比，火星大气极为稀薄。当前火星地表大气压从奥林匹斯山（Olympus Mons）上的 30 Pa（0.0044 psi）低值，到赫拉斯平原（Hellas Planitia）上的超过 1,155 Pa（0.1675 psi）高值不等，平均地表压力约为 600 Pa（0.087 psi）\(^{\text{[124]}}\)。火星上能够达到的最高大气密度，相当于地球表面上方 35 公里（22 英里）处的大气密度 \(^{\text{[125]}}\)。因此火星的平均地表压力仅为地球 101.3 kPa（14.69 psi）的 0.6\%。火星大气的尺度高度约为 10.8 公里（6.7 英里）\(^{\text{[126]}}\)，高于地球的 6 公里（3.7 英里），这是因为火星表面重力仅为地球的约 38\% \(^{\text{[127]}}\)。

火星大气由约 96\% 的二氧化碳、1.93\% 的氩气与 1.89\% 的氮气组成，并含有微量的氧气与水 \(^{\text{[2][128][122]}}\)。大气中含有约 1.5 μm 大小的尘埃颗粒，使得从地表看去火星天空呈黄褐色 \(^{\text{[129]}}\)；由于悬浮的铁氧化物颗粒，它有时也呈粉红色 \(^{\text{[22]}}\)。火星大气中的甲烷浓度从北半球冬季约 0.24 ppb 波动到夏季约 0.65 ppb \(^{\text{[130]}}\)。其寿命估计为 0.6 至 4 年 \(^{\text{[131][132]}}\)，因此其存在意味着必须有活跃的气体来源。甲烷最可能由非生物过程产生，例如含水、二氧化碳与富含橄榄石的岩石之间的蛇纹石化反应，而橄榄石在火星上广泛存在 \(^{\text{[133]}}\)。然而，甲烷也可能由火星生命产生 \(^{\text{[134]}}\)。
\begin{figure}[ht]
\centering
\includegraphics[width=8cm]{./figures/a9dbad45b69fb2ac.png}
\caption{由 MAVEN 在紫外波段观测到的火星逃逸大气（碳、氧与氢）。\(^{\text{[135]}}\)} \label{fig_Mars_16}
\end{figure}
与地球相比，火星大气中较高的二氧化碳浓度与较低的地表气压可能是声音在火星上衰减更为明显的原因；在火星上，自然声源除风以外十分稀少。基于“毅力号”（Perseverance）火星车获取的声学录音，研究者得出结论：当地声速在 240 Hz 以下的频率约为 240 m/s，而在更高频率下约为 250 m/s \(^{\text{[136][137]}}\)。

火星上已观测到极光现象 \(^{\text{[138][139][140]}}\)。由于火星缺乏全球磁场，其极光的类型与分布均不同于地球；与地球极光主要局限于极区不同，火星极光可遍布整颗行星 \(^{\text{[141][142]}}\)。2017 年 9 月，NASA 报道火星表面辐射水平在一次大规模且出乎意料的太阳风暴期间暂时增加至平时的两倍，并伴随出现亮度为此前观测到的 25 倍的极光 \(^{\text{[142][143]}}\)。
\subsubsection{气候}
\begin{figure}[ht]
\centering
\includegraphics[width=8cm]{./figures/8da4b831c9552fc8.png}
\caption{2001 年 7 月无全球尘暴的火星（左）与发生全球尘暴的火星（右），可见不同的水冰云层，由哈勃太空望远镜（Hubble Space Telescope）拍摄。} \label{fig_Mars_17}
\end{figure}
\begin{figure}[ht]
\centering
\includegraphics[width=8cm]{./figures/8546f9a4ee0dacc6.png}
\caption{火星北极（左）与南极（右）在北半球与南半球夏季之间二氧化碳冰（非水冰）覆盖变化的渲染图。} \label{fig_Mars_18}
\end{figure}
火星具有类似地球的季节变化，在其北半球与南半球之间交替出现。此外，相较于地球轨道，火星的轨道离心率更大：当南半球为夏季、北半球为冬季时，火星位于近日点；当南半球为冬季、北半球为夏季时，火星位于远日点。其结果是，南半球的季节变化更为极端，而北半球的季节则比本应出现的情况更为温和。南半球的夏季温度可比北半球对应季节高出多达 30 °C（54 °F）\(^{\text{[144]}}\)。

火星地表温度在赤道夏季的最高可达约 35 °C（95 °F），最低可达约 −110 °C（−166 °F）\(^{\text{[16]}}\)。这一极大的温差源于火星稀薄的大气无法储存多少太阳热量、大气压仅为地球大气压的约 1\%，以及火星土壤的低热惯量 \(^{\text{[145]}}\)。火星距离太阳的平均距离是地球的 1.52 倍，因此其接收到的阳光量仅为地球的 43\% \(^{\text{[146][147]}}\)。

火星拥有太阳系中最大的沙尘暴，其风速可超过 160 km/h（100 mph）。其规模可从局部区域的小型风暴到覆盖整颗行星的全球性沙尘暴。它们通常在火星最接近太阳时出现，并已被证明会提升全球气温 \(^{\text{[148]}}\)。季节变化还会使极冠形成干冰覆盖层 \(^{\text{[149]}}\)。
\subsection{水文}
\begin{figure}[ht]
\centering
\includegraphics[width=8cm]{./figures/8605fe891fd8f42e.png}
\caption{火星含有水，但主要以被尘埃覆盖的极地冰层形式存在，如本图所示。} \label{fig_Mars_19}
\end{figure}
尽管火星含有较大量的水，但其中大部分以被尘埃覆盖的水冰形式存在于火星极地冰冠之中 \(^{\text{[150][151][152][153][85]}}\)。如果南极冰冠中的水冰全部融化，其体积足以在火星大部分表面形成一层约 11 米（36 英尺）深的水层 \(^{\text{[154]}}\)。

由于火星的大气压不足地球的 1\%，液态水无法在其地表长期存在 \(^{\text{[155]}}\)。只有在最低洼的区域，气压与温度才可能在短时间内满足液态水存在的条件 \(^{\text{[47][156]}}\)。

虽然火星大气中的水含量很少，但仍足以形成水冰云，以及不同形式的雪与霜，这些现象常与二氧化碳干冰雪混合出现 \(^{\text{[157]}}\)。
\subsubsection{过去的水圈}
\begin{figure}[ht]
\centering
\includegraphics[width=8cm]{./figures/cc4f1bb7411d46a7.png}
\caption{艺术家笔下四十亿年前的火星想象图。} \label{fig_Mars_20}
\end{figure}
火星上可见的各种地貌强烈表明液态水曾存在于其地表。巨大的线性冲蚀带，称为溢流槽（outflow channels），在约 25 处横贯火星表面。一般认为，这些地貌记录了地下含水层中大量水突发释放所造成的侵蚀，但也有部分结构被假设可能由冰川或熔岩作用形成 \(^{\text{[158][159]}}\)。其中一处规模较大的例子——马阿迪姆谷（Ma'adim Vallis）长达 700 公里（430 英里），远大于地球上的科罗拉多大峡谷，其宽度达 20 公里（12 英里），部分区域深达 2 公里（1.2 英里）。它被认为是在火星早期由流动水切割而成 \(^{\text{[160]}}\)。这些溢流槽中最年轻的一部分被认为仅在几百万年前形成 \(^{\text{[161]}}\)。

在火星的其他地区，尤其是火星表面最古老的区域，遍布更为细密、树状分支状的谷网。这些谷地形态及其分布强烈暗示它们是由降水产生的径流在火星早期历史中切割而成。地下水流动与地下水侧蚀（groundwater sapping）可能在部分谷网中发挥次要作用，但几乎所有情况下，最根本的侵蚀来源应为降水 \(^{\text{[162]}}\)。

在陨石坑和峡谷壁上，有数千处与地球类似的沟槽（gullies）。这些沟槽多见于南半球高地，其方向朝向赤道，并全部位于纬度 30° 以极向一侧。许多研究者认为其形成过程涉及液态水，可能源自冰融化 \(^{\text{[163][164]}}\)，而另一些人则提出可能由二氧化碳霜或干尘移动形成 \(^{\text{[165][166]}}\)。尚未发现因风化而部分退化的沟槽，也未发现覆盖于其上的撞击坑，说明这些地貌非常年轻，并可能仍处于活动状态 \(^{\text{[164]}}\)。其他地质特征，如保存在陨石坑中的三角洲与冲积扇，则进一步支持火星早期某一或数个阶段出现过更温暖、更潮湿的环境 \(^{\text{[167]}}\)。此类环境必然要求火星表面大面积存在陨石坑湖泊，而矿物学、沉积学与地貌学的独立证据均支持这一点 \(^{\text{[168]}}\)。更多证据来自于特定矿物（如赤铁矿和针铁矿）的发现，这些矿物往往在水的参与下形成 \(^{\text{[169]}}\)。
\subsubsection{水的观测历史与相关证据的发现}
\begin{figure}[ht]
\centering
\includegraphics[width=8cm]{./figures/2007889d13f8e0f2.png}
\caption{如本图中这座名为柯罗廖夫（Korolev）的极地陨坑所示，火星表层的水冰在某些区域清晰可见，此为其三维投影图像。} \label{fig_Mars_21}
\end{figure}
2004 年，“机遇号”（Opportunity）探测器检测到黄钾铁矾（jarosite）这一矿物。它仅能在酸性水的存在下形成，因此表明火星上曾经存在过水 \(^{\text{[170][171]}}\)。2007 年，“勇气号”（Spirit）探测器发现高浓度二氧化硅沉积，显示火星过去具有潮湿环境；2011 年 12 月，“机遇号”再次在火星表面发现了石膏，这同样是一种在有水条件下形成的矿物 \(^{\text{[172][173][174]}}\)。估计火星上地幔上部所含的水量——以火星矿物中所含氢氧根离子的形式表示——在水含量为 50–300 ppm 的范围内，与地球相当甚至更高，此量足以在火星表面形成 200–1,000 米（660–3,280 英尺）深的全球性水层 \(^{\text{[175][176]}}\)。

2013 年 3 月 18 日，NASA 报告称“好奇号”（Curiosity）火星车的仪器在多份岩石样品中发现矿物水合现象，推测为水合硫酸钙。这些样品包括破裂的“Tintina”岩与“Sutton Inlier”岩碎片，以及诸如“Knorr”岩与“Wernicke”岩等其他岩石中的脉体与结核 \(^{\text{[177][178]}}\)。使用火星车 DAN 仪器的分析表明，地表以下 60 厘米（24 英寸）范围内存在地下水，其含水量高达 4\%；这些测量是在火星车从布拉德伯里着陆点（Bradbury Landing）驶向格伦尔格（Glenelg）地形中黄刀湾（Yellowknife Bay）区域的过程中获得的 \(^{\text{[177]}}\)。2015 年 9 月，NASA 宣布基于光谱仪对坡面暗色区域的分析，发现了关于重复坡线（recurring slope lineae）中水合盐卤（hydrated brine）流动的强有力证据 \(^{\text{[179][180][181]}}\)。这些条纹在火星夏季向坡下流动，此时温度高于 −23 °C；在更低温度下则冻结 \(^{\text{[182]}}\)。这些观测支持了早先基于其形成时间与增长速率提出的假说，即这些暗条纹源于地表下方水的流动 \(^{\text{[183]}}\)。然而后续研究提出，这些坡线可能是干燥的颗粒流动，其形成过程中水的作用至多有限 \(^{\text{[184]}}\)。目前关于火星地表液态水的存在、范围与作用仍缺乏定论 \(^{\text{[185][186]}}\)。

研究者推测，火星北部的低洼平原曾被一片深达数百米的海洋所覆盖，尽管这一理论仍存在争议 \(^{\text{[187]}}\)。2015 年 3 月，科学家提出这一古海洋可能与地球的北冰洋规模相当。这一结论基于当代火星大气中氕与氘的比值与地球的对应比值的比较。火星大气中的氘含量（D/H = 9.3 ± 1.7 × 10\(^{-4}\)）是地球（D/H = 1.56 × 10\(^{-4}\)）的五至七倍，这表明古代火星曾具有显著更高的水量。此前，“好奇号”（Curiosity）曾在盖尔陨坑（Gale Crater）检测到较高的氘比值，尽管不足以明确证明海洋曾经存在。然而，有科学家警告称这些结果尚未得到进一步确认，并指出当前的火星气候模型尚未显示火星过去具备足够温暖的条件以支持液态水体的长期存在 \(^{\text{[188]}}\)。在北极冰冠附近，火星快车号（Mars Express）发现宽 81.4 公里（50.6 英里）的柯罗廖夫陨坑（Korolev Crater）中约储存有 2,200 立方公里（530 立方英里）的水冰 \(^{\text{[189]}}\)。

2016 年 11 月，NASA 报告在乌托邦平原（Utopia Planitia）地区发现大量地下冰，其水量估计相当于苏必利尔湖（Lake Superior）的水量（12,100 立方公里 \(^{\text{[190]}}\)）\(^{\text{[191][192]}}\)。在 2018 至 2021 年期间，“ExoMars 痕量气体轨道器”（Trace Gas Orbiter）在水手谷（Valles Marineris）峡谷系统中探测到疑似水的迹象，可能为地下冰的存在 \(^{\text{[193]}}\)。
\subsection{轨道运动}
\begin{figure}[ht]
\centering
\includegraphics[width=8cm]{./figures/620ad9bc1db5b447.png}
\caption{火星及太阳系内侧行星的轨道。} \label{fig_Mars_22}
\end{figure}
火星与太阳的平均距离约为 2.3 亿公里（1.43 亿英里），其公转周期为 687 个地球日。火星的太阳日（sol）仅比地球日稍长，为 24 小时 39 分 35.244 秒 \(^{\text{[194]}}\)。一个火星年等于 1.8809 个地球年，即 1 年 320 天 18.2 小时 \(^{\text{[2]}}\)。从地球转移到火星所需跨越的引力势差及相应的变轨速度（delta-v）是所有行星间转移中仅次于地球—金星转移的第二低 \(^{\text{[195][196]}}\)。

火星自转轴相对其轨道平面的倾角为 25.19°，与地球的自转轴倾角十分相近 \(^{\text{[2]}}\)。因此火星亦存在类似地球的季节变化，只是由于火星轨道周期更长，其季节持续时间几乎是地球的两倍。在当今时代，火星北极的指向大致接近天津四（Deneb）\(^{\text{[21]}}\)。

火星的轨道偏心率约为 0.09，相对明显；在太阳系的其他七颗行星中，只有水星的轨道偏心率更大。已知火星过去的轨道要圆得多。例如约 135 万年前，火星的轨道偏心率仅约 0.002，远小于今日地球的偏心率 \(^{\text{[197]}}\)。火星轨道偏心率的变化周期约为 96,000 个地球年，相比之下地球为 100,000 年 \(^{\text{[198]}}\)。

火星与地球的最近接（冲日）发生在 779.94 天的会合周期中。需要注意，这与“合日”（conjunction）不同，后者是地球与火星分处太阳两侧呈一直线。火星两次冲日间的时间即为其会合周期，平均为 780 天；但实际间隔可能在 764 至 812 天之间变化 \(^{\text{[198]}}\)。由于两者轨道均为椭圆形，最近距离可在约 5,400 万至 1.03 亿公里（3,400 万至 6,400 万英里）之间变化，对应角直径也随之变化 \(^{\text{[199]}}\)。在最远时，火星与地球的距离可达 4.01 亿公里（2.49 亿英里）\(^{\text{[200]}}\)。火星每约 2.1 年从地球看会出现一次冲日。在 2003、2018 与 2035 年，火星在近日点附近发生冲日，而 2020 与 2033 年的冲日事件则特别接近于典型的近近日点冲日 \(^{\text{[201][202][203]}}\)。








